%%
% BIThesis 本科毕业设计论文模板 —— 使用 XeLaTeX 编译 The BIThesis Template for Undergraduate Thesis
% This file has no copyright assigned and is placed in the Public Domain.
%%

\begin{appendices}
  \section{主要符号说明}
  为便于阅读，本文使用的部分主要符号说明如下。$x,y$ 表示单体在世界坐标系下的位置；$\phi$ 表示单体航向角；$v$ 与 $\omega$ 分别表示单体线速度和角速度；$L_0$ 表示等效腿长；$\Phi_0$ 表示等效腿摆角；$\theta_i$ 表示链式系统中第 $i$ 个铰接角；$X$ 表示状态向量；$U$ 表示控制输入向量；$K$ 表示 LQR 反馈增益矩阵；$Q$ 与 $R$ 分别表示 LQR 性能指标中的状态权重矩阵和输入权重矩阵。

  \section{实验评价指标说明}
  链式自重组轮足机器人实验可采用以下指标进行评价：速度跟踪误差用于衡量单体对期望线速度和角速度的跟踪能力；姿态误差用于衡量机体俯仰角、滚转角和航向角的稳定性；铰接角误差用于衡量链式构型保持能力；控制输入峰值和控制输入变化率用于评价执行器负担和控制平滑性；稳定时间用于描述系统受到扰动或指令变化后的恢复速度。

\end{appendices}
